\chapter{EXPERIMENTAL DESIGN AND PRE-ANALYSIS}

Proper experimental planning is important for developing an efficient fluoride adsorption system. In the present work, GTMAC-modified coconut fiber is selected as the proposed adsorbent material for fluoride removal studies. Design of Experiments \gls{doe} using Central Composite Design \gls{ccd} is adopted for planning the adsorption experiments and optimization studies. This chapter presents the experimental specifications, pre-analysis work, adsorption models, and experimental techniques selected for the study.  
\section{Specifications for the Design}
The present study focuses on fluoride removal from aqueous solution using \gls{gtmac}-modified coconut fiber as the proposed adsorbent material. Batch adsorption studies are planned under controlled laboratory conditions by varying important adsorption parameters such as pH, adsorbent dose, fluoride concentration, and contact time. 
The selected operating ranges proposed for the adsorption studies are shown below. 
\begin{table}[h]
	\centering
	\caption{Experimental Parameters and Their Ranges}
	\label{tab:parameters}
	\begin{tabular}{|l|l|}
		\hline
		\textbf{Parameter} & \textbf{Range} \\
		\hline
		pH & 2 -- 7 \\
		Adsorbent Dose & 0.1 -- 1 g/100 mL \\
		Initial Fluoride Concentration & 2 -- 70 ppm \\
		Contact Time & 30 -- 180 min \\
		Adsorbent &\gls{gtmac}-modified coconut fiber \\
		Fluoride Measurement & Fluoride Ion-Selective Electrode \\
		Experimental Design & Central Composite Design \gls{ccd} \\
		\hline
	\end{tabular}
\end{table}

The adsorption experiments are proposed to be designed using Design-Expert software with approximately 30 CCD experimental runs for adsorption studies and optimization.
%\begin{enumerate}
%\item Specifications for the Design
%\item Pre analysis work for the design or Models used
%\item Design methodology in detail
%\item Design Equations
%\item Experimental techniques (if any)
%\end{enumerate}
%Apart from the aforementioned sections, you can add sections as per the requirements of the project in consultation with your guide.

\section{Pre-analysis Work and Models Used }
Literature analysis was carried out to study previous research related to fluoride adsorption, biomass-based adsorbents, surface modification techniques, and adsorption modelling. Coconut fiber was selected as the proposed adsorbent because of its low cost, eco-friendly nature, and easy availability. Since raw coconut fiber has limited adsorption efficiency, Glycidyl Trimethyl Ammonium Chloride \gls{gtmac} was selected as the proposed surface modifying agent to improve fluoride adsorption through electrostatic attraction. 

The adsorption behaviour is proposed to be analyzed using the Langmuir adsorption isotherm model, represented by: 
isotherm model, represented by:

\begin{equation}
	q_e = \frac{q_{max} b C_e}{1 + b C_e}
	\label{eq:langmuir2}
\end{equation}

where:
\begin{itemize}
	\item $q_e$ -- adsorption capacity at equilibrium
	\item $q_{max}$ -- maximum adsorption capacity
	\item $C_e$ -- equilibrium fluoride concentration
	\item $b$ -- Langmuir adsorption constant
\end{itemize}

Design of Experiments (DOE) using Central Composite Design (CCD) was selected for planning adsorption studies and optimization. The adsorption data generated from the experiments are proposed to be analyzed using Design-Expert software.

The percentage fluoride removal is proposed to be calculated using:

\begin{equation}
	\% \text{ Removal} = \frac{C_0 - C_e}{C_0} \times 100
	\label{eq:removal}
\end{equation}

The adsorption capacity is proposed to be calculated using:

\begin{equation}
	q_e = \frac{(C_0 - C_e)V}{m}
	\label{eq:capacity}
\end{equation}	

\section{Experimental techniques }
The following instruments and techniques were used during the experimental work:

\begin{table}[h]
	\centering
	\caption{Instruments and Techniques Used}
	\label{tab:instruments}
	\begin{tabular}{|l|l|}
		\hline
		\textbf{Experimental Stage} & \textbf{Technique / Instrument} \\
		\hline
		Biomaterial Preparation & Hot air oven, grinder, sieve \\
		Surface Modification & Magnetic stirrer, pH meter \\
		\gls{gtmac} Functionalization & Sonication and alkaline treatment \\
		Batch Adsorption Studies & Orbital shaker \\
		Fluoride Measurement & Fluoride Ion-Selective Electrode \\
		Surface Characterization & \gls{ftir} and \gls{sem} \\
		Statistical Analysis & Design-Expert Software \\
		\hline
	\end{tabular}
\end{table}

This chapter presented the experimental specifications, adsorption models, pre-analysis 
work, and experimental techniques proposed for fluoride adsorption studies using 
\gls{gtmac}-modified coconut fiber. The selected adsorption parameters, \gls{doe}-based 
experimental design, and proposed analytical techniques discussed in this chapter form 
the basis for the proposed methodology and implementation workflow discussed in the 
next chapter.

%\vspace{0.75cm}
%
% \textbf{The chapters should not end with figures, instead bring the paragraph explaining about the figure at the end followed by a summary paragraph.}
%
%After elaborating the various sections of the chapter (From Chapter 2 onwards), a summary paragraph should be written discussing the highlights of that particular chapter. This summary paragraph should not be numbered separately. This paragraph should connect the present chapter to the next chapter.  